\documentclass[11pt,letterpaper]{article}
\usepackage[margin=0.85in,headheight=15pt]{geometry}
\usepackage[T1]{fontenc}\usepackage{lmodern}
\usepackage{amsmath,booktabs,array,graphicx,xcolor,xurl,hyperref,fancyhdr}
\definecolor{ink}{HTML}{17334A}\definecolor{accent}{HTML}{187E91}
\hypersetup{colorlinks=true,linkcolor=ink,urlcolor=accent,pdftitle={Necromatcher: Fixed Row-Time Sensitivity},pdfauthor={UpstreamDrift Development}}
\pagestyle{fancy}\fancyhf{}\fancyhead[L]{\small\textcolor{ink}{\textbf{NECROMATCHER} / ROW-TIME SENSITIVITY}}
\fancyhead[R]{\small Technical Supplement}\fancyfoot[L]{\small Fixed-Motion Diagnostic / Physical Timing Unknown}\fancyfoot[R]{\thepage}
\setlength{\parindent}{0pt}\setlength{\parskip}{6pt}\setlength{\emergencystretch}{2em}\renewcommand{\arraystretch}{1.15}
\newcommand{\code}[1]{\nolinkurl{#1}}\newcommand{\hash}[1]{{\footnotesize\nolinkurl{#1}}}
\begin{document}
{\small\textcolor{accent}{UPSTREAMDRIFT / REPEATABLE RESEARCH METHODS}}\par
{\LARGE\bfseries\textcolor{ink}{Fixed Row-Time Sensitivity}}\par
{\large Tiger V17 and Hogan V14 / Actual Read-Only Diagnostic}\par
\section{Executed Scope and Fixed State}
The diagnostic assessed 56 preregistered cases on four saved forearm motions:
44 supported cases and 12 complete-role support rejections. It changed only an
explicit encoded-raster row-time assumption. Saved camera, native geometry,
body observations, authored shaft evidence and spline coefficients remained fixed.
No optimizer, new fit, library write, registration or readout estimator ran.

Tiger uses 1280-by-720 images and Hogan 320-by-240 images; their pixel errors must
be interpreted separately. Training fragments and previously seen evaluation
fragments retain their existing roles. This is not untouched validation, physical
sensor calibration or a test of every possible shutter model.

\section{Authored Row-Time Equation}
For endpoint $k$ of source fragment $i$, the diagnostic uses
\begin{equation}
t_{ik}=t_i^{\mathrm{PTS}}+d f\Delta_{\mathrm{enc}}
\left(\frac{y_{ik}}{H-1}-\frac12\right),
\qquad f\in\{0,\tfrac14,\tfrac12,1\},\quad d\in\{-1,+1\}.
\end{equation}
$H$ is the full encoded raster height. $\Delta_{\mathrm{enc}}$ is the uniform
source PTS frame interval, approximately 33.367 ms for Tiger and 33.333 ms for
Hogan. The nonnegative readout assumption is $f\Delta_{\mathrm{enc}}$; $d$ gives
scan direction. These units refer to the container clock, not verified physical
sensor time. The centered reference row and full-raster mapping are authored,
with no crop, field, stabilization, exposure or sensor-height mapping inferred.

Each endpoint is evaluated against the projected native infinite shaft axis at
its own time. For $M$ observed fragments, the raw diagnostic is
\begin{equation}
R(f,d)=\sqrt{\frac{1}{2M}\sum_{i=1}^{M}\sum_{k=1}^{2}e_{ik}(t_{ik})^2},
\qquad \Delta R=R(f,d)-R(0,+1).
\end{equation}
Here $e$ is signed perpendicular image distance in pixels. It is unweighted and
separate from body RMS or an optimization objective. Authored fragment endpoints
are not physical club endpoints. Nonzero row timing has no common straight
bearing angle because the two endpoint comparisons use different times; angle
diagnostics remain null for every nonzero supported case.

\textbf{Grid accounting:} one zero case with direction $+1$, plus both directions
at three nonzero fractions, gives seven cases per fit/role. Four fits and two roles
give 56 unique cases. No duplicate zero directions or missing values are counted
as additional evidence.
\newpage
\section{Seen Evaluation Sensitivity}
\begin{center}\includegraphics[width=.94\linewidth]{assets/seen-delta-rms.pdf}\end{center}
{\small\textbf{Figure 1.} Separate pixel scales for Tiger and Hogan. Markers are the
actual supported grid points; connecting segments aid reading and are not a fitted
readout curve. The two arms share neither a measured sensor readout nor a new camera.
Tiger training is omitted from the plot because every nonzero case is unsupported.}

\begin{center}\small\begin{tabular}{@{}llrrr@{}}\toprule
Evidence Role & Arm & Baseline RMS & Minimum $\Delta R$ & Maximum $\Delta R$\\\midrule
Tiger Seen & Control & 151.999 & -2.2256 & +2.0917 \\
Tiger Seen & Variant & 149.966 & -2.0679 & +1.9255 \\
Hogan Seen & Control & 30.976 & -0.0131 & +0.0136 \\
Hogan Seen & Variant & 30.930 & -0.0204 & +0.0209 \\
Hogan Training & Control & 28.187 & -0.0269 & +0.0269 \\
Hogan Training & Variant & 28.174 & -0.0251 & +0.0251 \\
\bottomrule\end{tabular}\end{center}
{\small Values are in pixels; displayed values are rounded. Complete precision,
endpoint times, errors, status and rejection reasons accompany all 56 JSON/CSV rows.}

Tiger seen errors change by at most about 2.23 pixels around roughly 150 pixels
of mismatch. Hogan training changes by less than 0.027 pixels; seen evaluation by
less than 0.021 pixels. This fixed scan does not resolve the historical club-axis
mismatch. It neither calibrates shutter timing nor rules out arbitrary physical
shutter, exposure, flexure or camera models outside these explicit assumptions.
\newpage
\section{Whole-Role Support and Exact Baselines}
Before any shifted motion evaluation, every observed endpoint time must lie inside
the saved spline interval. If one endpoint is outside, the entire evidence role
is rejected. No clamp, extrapolation, dropped endpoint, narrowed role or invented
zero error is allowed. Authored abstentions remain absent observations.

\begin{center}\begin{tabular}{@{}lrr>{\raggedright\arraybackslash}p{.39\linewidth}@{}}\toprule
Player / Role & Supported & Rejected & Reason or Scope\\\midrule
Tiger Training & 2 & 12 & Both arms' six nonzero scans cross boundary support.\\
Tiger Seen & 14 & 0 & Two arms, seven supported grid points each.\\
Hogan Training & 14 & 0 & Original interior training fragment remains supported.\\
Hogan Seen & 14 & 0 & Two arms, seven supported grid points each.\\\midrule
Total & 44 & 12 & 56 unique whole-role cases.\\\bottomrule
\end{tabular}\end{center}

All twelve rejection records retain the reason: ``Complete evidence role requires
all endpoint times inside spline support.'' Tiger training fragments lie at source
indices 0 and 209, the saved interval boundaries. Opposite scan signs shift an
endpoint beyond opposite support ends. Two zero training cases are supported;
this does not justify a nonzero Tiger training curve or a conclusion of zero sensitivity.
Interior source evidence is needed before testing those roles under signed shifts.

All eight zero cases reproduce unchanged legacy raw RMS and endpoint errors exactly.
The four training baselines also reproduce saved canonical shaft assessments.
The legacy zero angular errors are retained; all supported nonzero angular entries
remain null. Only the evidence actually observed contributes to the unweighted RMS.

\section{Strict Saved-Final Reconstruction}
Canonical preserved-spline loading and the public strict trajectory initializer
reconstruct each saved final motion without optimization or new feasible projection.
Knot times, free order and physical coefficient identities are checked exactly.
All dense source positions match saved samples at absolute tolerance $10^{-12}$,
with zero relative tolerance. Recomputed body RMS matches the saved value at absolute
tolerance $10^{-10}$, with zero relative tolerance; it is not a relative 25-pixel guard.

The read-only initializer's reference $q_0$ is a reconstruction input from the
saved motion. It is explicitly not claimed to be the original job's soft prior.
Body RMS, camera, geometry and coefficient identity are controls, not optimized
results of this scan. Canonical library binding rechecks actual capture, source
PNG/BGR identity, camera, evidence and full source-clock hashes before numerical use.
\newpage
\section{Provenance and Independent Preservation}
Execution producer is \hash{e0042ba61870f88543cb1305e3f6d21903c97fb8}.
The existing four saved motions remain the earlier 30-budget Tiger V17/Hogan V14
forearm candidates, not the partial 120-budget Tiger V18 results. All candidates
remain nonconverged and rejected; this diagnostic changes no stored scientific status.

\begin{tabular}{@{}p{.23\linewidth}p{.70\linewidth}@{}}\toprule
Record & SHA-256\\\midrule
Actual Assessment & \hash{750453ad2e2d9e02c0ada9eafcc9a18db3628aef9aa72e2fc9e794856d5dbe2f}\\
Independent Review & \hash{f3481f936e4ba619aa06ff0beee579eb3cd179e60b7c67b4595add8f5b59114c}\\
Finally Preservation & \hash{20ef0fdf8fc303cfcb1b7ca85aafb6dd5c918b3ccc5872961add8676088ab4c9}\\\bottomrule
\end{tabular}\par\smallskip
The independent review rehashed 14,767 source files and all 347 current library
files. Before/after source, runtime, whole-library, protocol and input identities
match exactly. No original bytes, index records or source images changed, and
there is no permitted library growth. The independent verifier reread saved results
and checked endpoint RMS/support records without repeating native execution.

\code{evidence/} preserves the exact assessment and both independent/finally receipts;
\code{input-pins.json} records their original paths, copy sizes and SHA-256.
\code{all-56-raw-rows.json} preserves complete case assessments, including unsupported
null entries. The CSV has one row per case with serialized endpoint arrays and blank
unsupported metrics; it does not silently substitute zeros. The vector scientific
plot and a rendered PNG are in \code{assets/}.

\section{What This Does Not Establish}
No readout parameter was fitted, no rolling-shutter cause identified, and no
physical sensor/exposure clock qualified. Full encoded raster is not known to
represent the sensor readout geometry. Video playback scale, sensor crop/field
mapping, motion blur, exposure integration, flexible shaft behavior, camera motion
and anatomical morphology remain uncertain. The scan does not measure club heads,
physical shaft endpoints, validated anatomy, clinical ROM or dynamics.

A source-only dense review now preserves 52 equally spaced images and root's visual
review, but no new points, phase/contact events or training/held-out roles have been
admitted. Tiger's broad moving fans and Hogan's raster scratch require abstention
and source-supported interpretation, not retrospective model-based labels. Future
timing evidence must preregister interior coverage and valid support margins.

The three-page partial solver-budget checkpoint and the earlier twelve-page
forearm report remain preserved. This report is a distinct actual fixed-motion
sensitivity record; it asserts no public deployment, issue closure, calibrated
capture or completed reconstruction. Numerical reproducibility and visual software
checks are separate from historical-motion qualification.
\end{document}
